% STATUS: DRAFT
\chapter{Takesaki duality}\label{ch:takesaki}

\begin{physmotiv}
Having added a clock to a type III algebra, we want to know what we got. The answer is the structural theorem of the first half of the book: the crossed product of a type III$_1$ algebra by its modular flow is a type II$_\infty$ factor, it has a \emph{trace}, and the original algebra is recovered as the fixed points of a dual symmetry that rescales the trace. The existence of the trace is a consequence of the KMS property of the state, so it comes ``for free'' from thermal physics. The trace is the thing that will define entropy.
\end{physmotiv}

\section{The dual action}

On $N=\M\rtimes_\sigma\R$ (frame 1, \cref{ch:crossed}) define
\begin{equation}
\theta_s=\Ad\,\ee^{\ii sp},\qquad s\in\R.
\end{equation}
Since $\ee^{\ii sp}x\ee^{-\ii sp}=x+s$ and $p$ commutes with $a\otimes\id$ and with $K$, we have $\theta_s(a\otimes\id)=a\otimes\id$ and $\theta_s(X)=X+s$. So $\theta$ is a one-parameter group of automorphisms of $N$ that fixes $\M$ and translates $X$. In terms of the elements $\int\dd u\,a(u)\ee^{\ii uX}$, it multiplies the mode $u$ by $\ee^{\ii us}$.

\begin{proposition}
The fixed-point algebra is $N^\theta=\M\otimes\id\cong\M$.
\end{proposition}
(Intuitively: only the mode $u=0$ is invariant.) So we recover $\M$ from $(N,\theta)$.

\section{The dual weight and the trace}

Define the conditional expectation-like map $T:N_+\to\M_+$ (an \emph{operator-valued weight}) by averaging over the dual action, $T(b)=\int\dd s\,\theta_s(b)$ suitably normalized; on $b=\int\dd u\,a(u)\ee^{\ii uX}$ it picks out the zero mode $a(0)$. The \emph{dual weight} is
\begin{equation}
\hat\varphi=\omega\circ T,\qquad\omega=\inner\Omega{\,\cdot\,\Omega}.
\end{equation}
It is a faithful normal semifinite weight on $N$, $\theta$-invariant, and for a function of $X$ alone, $\hat\varphi(f(X))=\int\dd X\,f(X)$ (normalization fixed by convention). Its modular flow is $\sigma^{\hat\varphi}_t=\Ad\ee^{-\ii tX}$ (it restricts to $\sigma^\omega_t$ on $\M$ and fixes $X$).

\begin{theorem}[trace on the crossed product]\label{thm:trace}
Let $\tau(b)=\hat\varphi\big(\ee^{X/2}\,b\,\ee^{X/2}\big)$ for $b\in N_+$. Then $\tau$ is a faithful normal semifinite \emph{trace} on $N$, and it is rescaled by the dual action:
\begin{equation}
\tau\circ\theta_s=\ee^{-s}\,\tau.
\end{equation}
Moreover, for any bounded Borel function $f$,
\begin{equation}
\tau\big(f(X)\big)=\int_{-\infty}^\infty\dd X\;\ee^{X}f(X).
\end{equation}
\end{theorem}

\begin{proofsketch}
\emph{Trace property.} Write $b=\int\dd u\,a(u)\ee^{\ii uX}$. The weight $\hat\varphi$ has modular flow $\Ad\ee^{-\ii tX}$, so $\hat\varphi_h=\hat\varphi(h^{1/2}\,\cdot\,h^{1/2})$ has modular flow $\Ad(h^{\ii t})\circ\Ad\ee^{-\ii tX}$ (Connes cocycle, \cref{ch:modham}). This is trivial iff $h=\ee^{X}$, and a weight with trivial modular flow is a trace. Direct check of the same fact: with $b e^{X/2}=\int\dd u\,a(u)\ee^{(\ii u+1/2)X}$ one finds $\tau(b\ad b)=\int\dd u\,\omega\big(\tilde a(u)\ad\tilde a(u)\big)$ with $\tilde a(u)=a(u-\ii/2)$ and $\tau(bb\ad)=\int\dd u\,\omega\big(\tilde a(u)\,\sigma_{-\ii}(\tilde a(u)\ad)\big)$, and these agree exactly by the KMS condition $\omega(c\,\sigma_{-\ii}(d))=\omega(dc)$ (\cref{ch:tomita}). \emph{Scaling:} since $\hat\varphi$ is $\theta$-invariant and $\theta_s(\ee^{X/2})=\ee^{s/2}\ee^{X/2}$, $\tau\circ\theta_s=\ee^{-s}\tau$. The last formula follows from $\hat\varphi(f(X))=\int f\,\dd X$.
\end{proofsketch}

So the \emph{KMS property of $\omega$ for its own modular flow is what makes $\tau$ tracial}. This is the mathematical heart of the whole book.

\begin{whatwrong}
Conventions matter and are easily misread across papers: (i) which sign of $\ii$ in $\Delta^{\pm\ii t}$; (ii) whether the dual action is $\Ad\ee^{\ii sp}$ or its inverse, i.e.\ whether the trace scales as $\ee^{-s}$ or $\ee^{+s}$; (iii) whether the variable called $x$ is $X$ or $-X$. We have fixed them in \cref{ch:crossed} (see also the style guide) and checked them in finite dimensions below; the physics statements (existence of the trace, finiteness for $x\le0$) are convention independent.
\end{whatwrong}

\section{Type of $N$ and Takesaki duality}

\begin{theorem}[Connes--Takesaki, as stated in \cref{ch:connes_class}]
If $\M$ is a $\sigma$-finite type III$_1$ factor, then $N=\M\rtimes_\sigma\R$ is a type II$_\infty$ factor with trace $\tau$ above. In general (type III, $\sigma$-finite), $N$ is semifinite: it has a f.n.s.\ trace.
\end{theorem}
The factor property uses that $\sigma$ has no ``periodic part'' (the centre of $N$ is the \emph{flow of weights}, which is trivial for III$_1$). Also:
\begin{theorem}[Takesaki duality]
$N\rtimes_\theta\R\;\cong\;\M\bar\otimes\Bnd{L^2(\R)}$.
\end{theorem}
So crossing twice returns the original algebra, up to tensoring with type I$_\infty$: the type is unchanged. Since $\M\otimes\Bnd{L^2}$ and $\M$ are Morita equivalent, one can say $\M$ is recovered from $(N,\tau,\theta)$: \emph{a type III$_1$ factor is the same data as a II$_\infty$ factor with a trace-scaling flow}.

\subsection*{II$_\infty$ versus II$_1$}
Because $\tau\circ\theta_s=\ee^{-s}\tau$, we have $\tau(\id)=\infty$ (indeed $\int\ee^{X}\dd X=\infty$). A type II$_\infty$ algebra has finite-trace \emph{projections}, and one can restrict to a corner. For $P=\chi(X\le0)\in N$,
\begin{equation}
\tau(P)=\int_{-\infty}^0\ee^{X}\dd X=1.
\end{equation}
The corner $PNP$ is a type II$_1$ factor with normalized trace $\tau(P\,\cdot\,P)$. In gravity the condition $X\le0$ will be the statement that the observer's energy is non-negative (\cref{ch:add_observer}); the maximum-entropy state will then have density matrix $P$ (\cref{ch:entropy_II,ch:clpw}).

\section{Finite-dimensional check}
For $\M=M_n$ (\cref{ch:crossed}), $N=M_n\otimes\Linf(\R_y)$ with centre variable $y=X+\log\rho$. A direct computation gives
\begin{equation}
\tau(F)=\int\dd y\,\ee^{y}\,\Tr F(y),\qquad F(y)\in M_n,
\end{equation}
and for $F=f(X)$, whose eigenvalues are $f(y-\log p_i)$ (since $X=y-\log\rho$),
\begin{equation}
\tau\big(f(X)\big)=\sum_i\int\dd y\,\ee^{y}f(y-\log p_i)=\sum_ip_i\int\dd X\,\ee^{X}f(X)=\int\dd X\,\ee^Xf(X),
\end{equation}
using $y=X+\log p_i$ and $\sum_ip_i=1$. This agrees with \cref{thm:trace}, and $\tau\circ\theta_s=\ee^{-s}\tau$ follows from $\theta_s:y\mapsto y+s$. Of course type I algebras already have traces; the point is that the formula $\tau(f(X))=\int\ee^Xf(X)\dd X$ is the one that survives when $\rho$ ceases to exist.

\begin{keyidea}
Cross the type III$_1$ factor by its modular flow: the result is a type II$_\infty$ factor with trace $\tau(f(X))=\int\ee^{X}f(X)\dd X$ coming from the KMS property. A dual flow $\theta_s$ rescales the trace by $\ee^{-s}$ and its fixed points are the original algebra. Cutting down to $X\le0$ gives a type II$_1$ corner.
\end{keyidea}

\section*{Exercises}
\begin{exercise}
Show that the fixed points of $\theta_s=\Ad\ee^{\ii sp}$ on $\{a\otimes\id,K+x\}''$ are $\M\otimes\id$: argue with the Fourier decomposition $b=\int\dd u\,a(u)\ee^{\ii uX}$.
\end{exercise}
\begin{exercise}
Verify the finite-dimensional trace formula. Use $N=M_n\otimes\Linf(\R_y)$ with $y=X+\log\rho$ and check (a) $\tau(FG)=\tau(GF)$ for $F,G\in N$ (trivial: $\Tr$ cyclic), (b) $\tau\circ\theta_s=\ee^{-s}\tau$ with $\theta_s:y\mapsto y+s$, (c) $\tau(f(X))=\int\ee^Xf(X)\dd X$ using $\sum p_i=1$.
\end{exercise}
\begin{exercise}
Show that $\tau(\chi(X\le0))=1$ and $\tau(\chi(X\ge0))=\infty$. Compute $\tau(\chi(X\le c))$ and show that all values in $(0,\infty)$ occur, so $N$ has projections of every dimension (type II).
\end{exercise}
\begin{exercise}
Verify the claim in the proof sketch: a weight $\varphi_h=\varphi(h^{1/2}\cdot h^{1/2})$ has $\sigma^{\varphi_h}_t=\Ad(h^{it})\circ\sigma^\varphi_t$ in finite dimensions ($\varphi=\Tr(\rho\,\cdot)$, $\varphi_h=\Tr(h^{1/2}\rho h^{1/2}\,\cdot)$ with $h$ commuting). Why is this trivial iff $h=\rho^{-1}$?
\end{exercise}
